\chapter{GIỚI THIỆU VỀ MỘT SỐ KHÁI NIỆM}

Chương 1. GIỚI THIỆU VỀ MỘT SỐ KHÁI NIỆM

\section{Tổng quan về thiết kế kiến trúc phần mềm}

\subsection{Khái niệm về thiết kế kiến trúc phần mềm}

Kiến trúc phần mềm (Software Architecture) là cấu trúc tổng thể của một hệ thống phần mềm, bao gồm các thành phần (components), mối quan hệ giữa chúng và các nguyên tắc chi phối việc thiết kế và phát triển hệ thống theo thời gian. Kiến trúc phần mềm không chỉ đơn thuần là bản vẽ kỹ thuật mà còn là tập hợp các quyết định thiết kế quan trọng, có ảnh hưởng sâu rộng đến toàn bộ vòng đời của sản phẩm từ giai đoạn phát triển, kiểm thử, triển khai cho đến bảo trì và nâng cấp. Theo định nghĩa của Viện Kỹ nghệ Phần mềm (SEI – Software Engineering Institute), kiến trúc phần mềm của một hệ thống là tập hợp các cấu trúc cần thiết để lập luận về hệ thống đó, bao gồm các thành phần phần mềm, các thuộc tính có thể quan sát được từ bên ngoài của các thành phần đó và mối quan hệ giữa chúng. Nói cách khác, kiến trúc phần mềm trả lời cho câu hỏi: “Hệ thống được tổ chức như thế nào--” ở mức độ trừu tượng cao nhất. Trong thực tế phát triển phần mềm hiện đại, kiến trúc phần mềm đóng vai trò như một bản thiết kế chi tiết (blueprint) cho cả hệ thống. Nó định nghĩa cách các module được phân tách và tổ chức, cách dữ liệu luân chuyển giữa các thành phần và cách hệ thống tương tác với các hệ thống bên ngoài. Một kiến trúc tốt cho phép nhiều nhóm phát triển làm việc song song trên các thành phần khác nhau mà không gây xung đột, đồng thời đảm bảo rằng khi ráp nối lại, toàn bộ hệ thống vận hành trơn tru như một thể thống nhất.

\subsection{Tại sao cần phải thiết kế kiến trúc phần mềm}

Thiết kế kiến trúc phần mềm không phải là một bước tùy chọn mà là một yêu cầu bắt buộc đối với mọi dự án phần mềm có quy mô từ trung bình trở lên. Những lý do chính bao gồm: Thứ nhất, đảm bảo tính module hóa và khả năng bảo trì. Khi hệ thống được phân tách thành các module độc lập với trách nhiệm rõ ràng, việc sửa lỗi hay nâng cấp một module không gây ảnh hưởng dây chuyền đến các phần khác. Điều này đặc biệt quan trọng đối với các hệ thống có vòng đời dài, nơi mã nguồn được duy trì bởi nhiều thế hệ lập trình viên khác nhau theo thời gian.

Thứ hai, tạo điều kiện cho phát triển song song. Một kiến trúc được phân tách rõ ràng cho phép nhiều nhóm cùng làm việc trên các thành phần khác nhau một cách độc lập. Trong dự án của nhóm em, ba dịch vụ N1, N2 và N3 được phát triển bởi ba nhóm khác nhau, mỗi nhóm chỉ cần tuân thủ API contract đã thống nhất mà không cần quan tâm đến chi tiết triển khai nội bộ của nhóm khác. Thứ ba, đáp ứng các yêu cầu phi chức năng. Các thuộc tính chất lượng như hiệu năng (performance), khả năng mở rộng (scalability), tính sẵn sàng (availability) và bảo mật (security) không thể được “thêm vào” sau khi hệ thống đã hoàn thiện. Chúng phải được thiết kế ngay từ đầu trong kiến trúc tổng thể. Ví dụ, quyết định sử dụng kiến trúc microservices cho hệ thống quản lý thư viện xuất phát từ yêu cầu mỗi dịch vụ phải có khả năng mở rộng độc lập theo tải trọng riêng. Thứ tư, giảm thiểu rủi ro kỹ thuật. Một kiến trúc được thiết kế cẩn thận giúp phát hiện sớm các vấn đề tiềm ẩn như nút thắt cổ chai về hiệu năng, xung đột dữ liệu hay phụ thuộc vòng (circular dependency) trước khi chúng trở thành những lỗi nghiêm trọng khó sửa trong giai đoạn triển khai. Thứ năm, tối ưu chi phí dài hạn. Mặc dù đầu tư thời gian cho thiết kế kiến trúc có thể làm chậm giai đoạn khởi động dự án, nhưng về lâu dài, nó giúp tiết kiệm đáng kể chi phí sửa lỗi, tái cấu trúc và mở rộng hệ thống. Các nghiên cứu trong ngành công nghiệp phần mềm đã chỉ ra rằng chi phí sửa một lỗi kiến trúc trong giai đoạn vận hành cao hơn gấp 10 đến 100 lần so với việc phát hiện và sửa nó trong giai đoạn thiết kế.

\section{Các kiểu kiến trúc thường dùng}

Trong lĩnh vực phát triển phần mềm, có nhiều kiểu kiến trúc (architectural styles) khác nhau, mỗi kiểu phù hợp với một lớp bài toán và bối cảnh cụ thể. Dưới đây là năm kiểu kiến trúc phổ biến được sử dụng rộng rãi trong công nghiệp phần mềm hiện đại.

\subsection{Kiến trúc phân lớp (Layered Architecture)}

Kiến trúc phân lớp, còn được biết đến với tên gọi kiến trúc n-tier, là một trong những kiểu kiến trúc lâu đời nhất và được áp dụng rộng rãi nhất trong phát triển phần mềm doanh nghiệp. Trong kiến trúc này, hệ thống được tổ chức thành các tầng (layer) xếp chồng lên nhau, mỗi tầng đảm nhiệm một trách nhiệm cụ thể và chỉ giao tiếp với tầng liền kề. Mô hình phổ biến nhất là kiến trúc ba tầng (three-tier) bao gồm: tầng trình bày (Presentation Layer) chịu trách nhiệm hiển thị giao diện và tiếp nhận tương tác từ người

dùng, tầng nghiệp vụ (Business Logic Layer) xử lý các quy tắc và logic nghiệp vụ của ứng dụng và tầng dữ liệu (Data Access Layer) đảm nhiệm việc đọc/ghi dữ liệu từ cơ sở dữ liệu hoặc các nguồn lưu trữ khác. Ưu điểm chính của kiến trúc phân lớp là tính đơn giản, dễ hiểu và dễ triển khai. Các lập trình viên có thể dễ dàng nắm bắt cấu trúc hệ thống và làm việc trên từng tầng một cách độc lập. Tuy nhiên, nhược điểm của kiến trúc này là khi hệ thống phát triển lớn, các tầng có xu hướng trở nên cồng kềnh (monolithic) và khó mở rộng độc lập. Ngoài ra, việc thay đổi một tầng thường kéo theo thay đổi ở các tầng liên quan do sự phụ thuộc chặt chẽ giữa các tầng.

\subsection{Kiến trúc Microservices}

Kiến trúc microservices là một phong cách kiến trúc trong đó ứng dụng được xây dựng như một tập hợp các dịch vụ nhỏ, độc lập, mỗi dịch vụ chạy trong tiến trình riêng và giao tiếp với nhau thông qua các giao thức nhẹ như HTTP/REST hoặc message queue. Mỗi microservice tập trung vào một chức năng nghiệp vụ cụ thể (single business capability) và có thể được phát triển, triển khai và mở rộng một cách độc lập. Trong dự án hệ thống quản lý thư viện, nhóm em đã áp dụng kiến trúc microservices với ba dịch vụ chính: Catalog Service (N1) quản lý danh mục sách, Circulation Service (N2) xử lý nghiệp vụ mượn trả và Identity Service (N3) đảm nhiệm xác thực người dùng. Mỗi dịch vụ có cơ sở dữ liệu riêng, được phát triển và triển khai độc lập bởi các nhóm khác nhau. Ưu điểm của kiến trúc microservices bao gồm khả năng mở rộng độc lập theo tải trọng của từng dịch vụ, sự cô lập lỗi (một dịch vụ gặp sự cố không làm sập toàn bộ hệ thống), khả năng sử dụng công nghệ khác nhau cho từng dịch vụ và hỗ trợ phát triển song song hiệu quả. Tuy nhiên, kiến trúc này cũng đặt ra những thách thức như độ phức tạp trong việc quản lý giao tiếp liên dịch vụ, tính nhất quán dữ liệu phân tán và chi phí vận hành hệ thống phân tán.

\subsection{Kiến trúc MVC (Model-View-Controller)}

MVC là một mẫu kiến trúc được sử dụng rộng rãi trong phát triển ứng dụng web, chia ứng dụng thành ba thành phần chính: Model (quản lý dữ liệu và logic nghiệp vụ), View (hiển thị giao diện người dùng) và Controller (tiếp nhận và xử lý yêu cầu từ người dùng, điều phối giữa Model và View).

Trong kiến trúc MVC, Controller đóng vai trò trung gian: khi người dùng tương tác với View (ví dụ: nhấn nút “Mượn sách”), Controller tiếp nhận yêu cầu, tương tác với Model để thực hiện logic nghiệp vụ (kiểm tra số lượng sách, tạo phiếu mượn) và trả về View cập nhật để hiển thị kết quả cho người dùng. Sự phân tách rõ ràng giữa ba thành phần giúp mã nguồn dễ bảo trì, dễ kiểm thử và cho phép phát triển song song giữa lập trình viên frontend (View) và backend (Model, Controller). Trong dự án N2, mặc dù backend được xây dựng theo kiến trúc REST API (không có View phía server), nhưng frontend Vue 3 của nhóm vẫn tuân theo tinh thần của MVC ở phía client: các component Vue đóng vai trò View, Vue Router kết hợp Pinia store đóng vai trò Controller và các service API đóng vai trò Model.

\subsection{Kiến trúc REST API}

REST (Representational State Transfer) là một phong cách kiến trúc được Roy Fielding đề xuất năm 2000, định nghĩa một tập hợp các ràng buộc để thiết kế các dịch vụ web. REST API tuân thủ các nguyên tắc: giao tiếp không trạng thái (stateless) – mỗi yêu cầu từ client đến server phải chứa đầy đủ thông tin cần thiết để xử lý, sử dụng các phương thức HTTP chuẩn (GET, POST, PUT, DELETE) để thao tác với tài nguyên được định danh bởi URI và dữ liệu được biểu diễn dưới các định dạng chuẩn như JSON hoặc XML. Dịch vụ N2 của nhóm em được thiết kế hoàn toàn theo phong cách REST API. Các tài nguyên như phiếu mượn (borrow transactions), phí phạt (fines), doanh thu (revenue) và đánh giá sách (reviews) đều được truy xuất thông qua các endpoint REST chuẩn với định dạng JSON. Ví dụ, endpoint GET /api/circulation/transactions trả về danh sách phiếu mượn và endpoint POST /api/circulation/borrow tạo một phiếu mượn mới. Kiến trúc REST API mang lại nhiều lợi ích: tính đơn giản và dễ hiểu, khả năng tương tác cao giữa các hệ thống khác nhau (vì sử dụng HTTP là giao thức phổ biến), khả năng mở rộng tốt nhờ tính không trạng thái và hỗ trợ caching ở nhiều cấp độ. Chính những ưu điểm này khiến REST API trở thành lựa chọn hàng đầu cho giao tiếp giữa các microservice trong các hệ thống phân tán hiện đại.

\subsection{Kiến trúc hướng sự kiện (Event-Driven Architecture)}

Kiến trúc hướng sự kiện là một phong cách kiến trúc trong đó các thành phần của hệ thống giao tiếp với nhau thông qua việc phát sinh (produce), phát hiện (detect) và phản

ứng (react) với các sự kiện. Một sự kiện là một thay đổi trạng thái đáng chú ý trong hệ thống, ví dụ như “sách đã được mượn”, “phiếu mượn đã quá hạn” hay “độc giả đã thanh toán phí phạt”. Thay vì gọi trực tiếp đến dịch vụ khác, các thành phần công bố sự kiện lên một kênh trung gian (event bus hoặc message broker), và các thành phần quan tâm sẽ đăng ký lắng nghe và phản ứng tương ứng. Trong dự án N2, nhóm em đã triển khai cơ chế PublishedEventLogs – một bảng ghi nhật ký sự kiện trong cơ sở dữ liệu, cho phép các dịch vụ khác (đặc biệt là N3 – Identity Service) có thể truy vấn và đồng bộ trạng thái dựa trên các sự kiện được N2 phát ra. Các loại sự kiện bao gồm: borrow.requested (yêu cầu mượn sách), borrow.approved (duyệt mượn), book.returned (sách đã trả), reader.notification (thông báo đến độc giả), cùng các sự kiện về thay đổi chính sách như BorrowPolicyUpdated và PricePolicyUpdated. Ưu điểm của kiến trúc hướng sự kiện là khả năng ghép nối lỏng (loose coupling) giữa các dịch vụ, khả năng mở rộng linh hoạt (có thể thêm consumer mới mà không ảnh hưởng đến producer) và phù hợp với các hệ thống cần xử lý luồng nghiệp vụ phức tạp, không đồng bộ. Đây cũng là nền tảng cho các kiến trúc tiên tiến hơn như CQRS (Command Query Responsibility Segregation) và Event Sourcing.

\section{So sánh các kiểu kiến trúc}

Mỗi kiểu kiến trúc có những ưu điểm và hạn chế riêng, phù hợp với các bối cảnh ứng dụng khác nhau. Bảng dưới đây tổng hợp và so sánh năm kiểu kiến trúc đã trình bày dựa trên các tiêu chí quan trọng trong phát triển phần mềm.

Bảng 1.1 So sánh các kiểu kiến trúc phần mềm

Tiêu chí Layered Microservices MVC REST API Event- (n-tier) Driven Độ phức Thấp Rất cao Trung bình Thấp – Cao tạp Trung bình Khả năng Hạn chế Rất tốt Trung bình Tốt Rất tốt mở rộng (scale cả (scale từng (stateless) khối) service) Khả năng Tốt với dự Rất tốt Tốt Tốt Tốt (loose bảo trì án nhỏ (module hóa coupling) cao) Hiệu năng Tốt (gọi nội Trung bình Tốt Tốt (caching Trung bình bộ) (network được) (async overhead) latency) Triển khai Đơn giản (1 Phức tạp Đơn giản Đơn giản Phức tạp artifact) (nhiều service) Phù hợp Ứng dụng Hệ thống Ứng dụng API public, Hệ thống với doanh lớn, nhiều web truyền mobile phân tán, nghiệp nhỏ – team thống backend real-time vừa Ví dụ thực Hệ thống Netflix, Ứng dụng Twitter API, Hệ thống tế ERP nội bộ Amazon, Spring GitHub API giao dịch Uber MVC, ngân hàng, Laravel IoT

\section{Lựa chọn kiến trúc cho dự án}

Từ bảng so sánh trên, có thể thấy không có một kiểu kiến trúc nào là “tốt nhất” cho mọi trường hợp. Việc lựa chọn kiến trúc phụ thuộc vào đặc thù của bài toán, quy mô nhóm phát triển và các yêu cầu phi chức năng. Đối với hệ thống quản lý thư viện số, nhóm em đã lựa chọn kết hợp nhiều kiểu kiến trúc dựa trên các lập luận sau: Microservices cho cấp độ hệ thống: Bài toán được phân công cho ba nhóm sinh

viên phát triển song song, mỗi nhóm có công nghệ và tiến độ riêng. Kiến trúc microservices cho phép mỗi nhóm làm việc độc lập trên service của mình, chỉ cần thống nhất API Contract. Đây là lựa chọn tối ưu về mặt tổ chức và phù hợp với yêu cầu của đề bài. Ngoài ra, trong thực tế, các thành phần của thư viện có tải trọng rất khác nhau — dịch vụ tìm kiếm sách (N1) có thể được truy vấn hàng nghìn lần mỗi giây trong khi dịch vụ mượn trả (N2) có tần suất thấp hơn. Microservices cho phép scale độc lập từng thành phần, tiết kiệm tài nguyên. REST API cho giao tiếp liên dịch vụ: Giao tiếp đồng bộ qua REST API được chọn cho các tương tác cần phản hồi ngay lập tức như kiểm tra số lượng sách trước khi cho mượn hoặc xác thực token. REST API có ưu điểm là đơn giản, dễ debug, được hỗ trợ rộng rãi và phù hợp với trình độ của nhóm. Tuy nhiên, nhóm cũng nhận thức được hạn chế của REST trong các tình huống cần giao tiếp bất đồng bộ hoặc khi một service bị chậm sẽ kéo theo timeout dây chuyền. Event-Driven cho đồng bộ trạng thái: Cơ chế Published Event Logs được áp dụng để giải quyết bài toán đồng bộ dữ liệu giữa các service mà không tạo ra phụ thuộc chặt chẽ. Khi N2 tạo một phiếu mượn, nó không gọi trực tiếp đến N3 để cập nhật báo cáo, mà ghi sự kiện vào bảng log. N3 có thể truy vấn log này theo chu kỳ để cập nhật dashboard. Cách tiếp cận này giảm coupling và cho phép N3 hoạt động độc lập ngay cả khi N2 gặp sự cố. Layered Architecture cho nội bộ từng service: Bên trong mỗi microservice, nhóm áp dụng kiến trúc phân tầng (Controllers → Business Logic → Data Access) để tổ chức mã nguồn rõ ràng, dễ hiểu và dễ bảo trì. Đây là mẫu kiến trúc được ASP.NET Core hỗ trợ sẵn thông qua cơ chế Dependency Injection và middleware pipeline. MVC cho frontend: Ở phía client, Vue 3 với Vue Router và Pinia store triển khai mô hình MVC một cách tự nhiên: View là các component Vue, Controller là Router + Store điều phối luồng dữ liệu và Model là các API service layer gọi đến backend. Sự kết hợp đa kiến trúc này — còn được gọi là hybrid architecture — phản ánh xu hướng thực tế trong công nghiệp phần mềm, nơi các hệ thống lớn hiếm khi tuân theo một kiểu kiến trúc thuần túy mà thường là sự pha trộn của nhiều phong cách khác nhau, mỗi phong cách giải quyết một khía cạnh cụ thể của bài toán.
